\chapter{Background Chapter}
\label{background}

\section{Introduction}

This chapter reviews static, dynamic, and hybrid approaches to deadlock detection, their limitations, and the motivation for reporting deadlocks as exceptions.

\section{Static Analysis}

Static analysis examines source code without executing it. Deadlock analyses typically look for cycles in resource-acquisition dependencies between threads.
These cycles indicate possible deadlocks, but some cannot occur during execution, leading to false positives \citep{dawson,chand,vivek,praun,cormac,williams}.

\citet{marino} analyze AJ, a Java extension whose atomic sets group class fields to prevent data races and atomicity violations by construction.
The analysis infers a partial order over these sets consistent with lock acquisition. Establishing this order proves deadlock freedom; otherwise, the tool reports potential deadlocks.
Implemented in the AJ-to-Java compiler, it uses synchronization annotations and demonstrated deadlock freedom for most evaluated programs without programmer intervention.
However, AJ's status as a research language limited the available subjects and their representativeness of practical concurrent programs.

\section{Dynamic Analysis}

Dynamic analysis examines execution traces to identify potential deadlock cycles. Observed behavior can improve precision, but coverage depends on the executions explored: a run without detected deadlocks does not establish deadlock freedom.
Instrumentation and analysis also introduce runtime and memory costs that can limit scalability.

\citet{mcsdk} present MulticoreSDK, a Java detector with two offline passes over an execution trace.
The first builds a graph of lock-event locations to identify potentially relevant locks; the second constructs a filtered graph using lock identities and reports its cycles.
Bytecode instrumentation records synchronization events for monitor instructions, synchronized methods, and the \emph{Lock} interface \citep{tanter}.
Compared with a traditional approach \citep{contest}, the reported evaluation found small analysis-time gains and roughly half the memory use for small applications; for large applications, the gain was about 33\%, with roughly one-tenth the memory use.

\citet{magicfuzzer} present MagicFuzzer, which records thread and lock events, extracts potential cycles from thread-specific lock dependencies, and re-executes the program to trigger them.
An active random scheduler explores executions until the candidate deadlocks are observed or a retry limit is reached; a proof of its claimed improvement in detection likelihood was left for future work.
The implementation uses Pin 2.9 \citep{pin} for C/C++ programs using Pthreads on Linux.
Evaluation on programs including MySQL \citep{mysql}, Firefox \citep{firefox}, and Chromium \citep{chromium} reported efficient execution and low memory use relative to tools such as MulticoreSDK \citep{mcsdk}.
However, workloads ranged from starting and closing an application to reproducing reported bugs, limiting the behavior explored.

\citet{sammati} present \emph{Sammati}, which intercepts Pthreads operations to detect deadlocks and recover without source changes.
It delays publication of critical-section memory updates until the enclosing locks are released. On detecting a deadlock, it rolls one participant back to the offending lock acquisition and discards its unpublished updates.
Memory is privatized by page rather than by instrumenting individual loads and stores. Following \citet{berger}, threads become \emph{cords}, single-threaded processes with separate address spaces that implement the shared-memory abstraction.
Lock-owner and waiting relationships support detection in $\mathcal{O}(n)$ time, where $n$ is the number of threads.
Recovery nevertheless incurs protection and privatization costs and is limited by effects that cannot safely be rolled back, including I/O, semaphore operations, condition signals, and shared-library state \citep{pyla}.
Serenity addresses some of these limitations through hybrid analysis, discussed below.

\citet{rx} propose Rx, which restores a server checkpoint after a failure and retries under a modified execution environment.
For timing-dependent failures such as deadlocks, a different execution may avoid the original failure.
A control unit selects environmental changes; memory-management wrappers apply them during recovery and otherwise forward calls to the standard libraries.
Changes are disabled once execution passes the failing region. The evaluation reported recovery in about 0.16 seconds, faster than restarting the web application.
Rx therefore emphasizes recovering service through re-execution rather than diagnosing the underlying bug.

\section{Mixing Static and Dynamic Analysis}

Hybrid approaches use static analysis to guide runtime checks, reducing the work required during execution.
Their effectiveness depends on the precision of the static model and the cost of the remaining dynamic checks.

\citet{grechanik} propose \emph{REDACT}, which identifies potential hold-and-wait cycles among SQL resources statically and uses a runtime controller to prevent them.
Developers first supply the application's database transactions. Instrumented interceptors then route SQL statements through the controller, which consults the precomputed cycles and delays conflicting statements while others proceed.
This reduces runtime analysis, but conservative cycles can cause unnecessary delays: some are infeasible or would already be prevented by the database engine.
False positives therefore add overhead and reduce parallelism.

\citet{serenity} present Serenity, a successor to Sammati \citep{sammati} that combines static analysis and selective instrumentation with runtime recovery.
It infers lock scopes to guide memory isolation, serializes disk I/O within critical sections, and tracks I/O and thread creation or termination operations.
When recovering, it selects a victim that has not performed I/O while holding locks.
Limitations remain for arbitrary I/O in restartable critical sections \citep{pyla}. Developers must identify lock operations hidden in functions or classes, and applications and their dependencies must be recompiled for instrumentation.

\section{Deadlocks as Exceptions}

The approaches above trade precision, coverage, runtime cost, and recovery capability. Automatic recovery can be restricted by critical-section side effects, while techniques designed for databases or servers may not transfer readily to other applications.

Concurrent Pascal \citep{hansen} addressed deadlocks at compile time by restricting the language, including a fixed number of processes and no recursion.
These constraints allowed analysis of a partial order of process interactions to establish properties such as deadlock freedom.

Runtime reporting also exists in restricted settings. Database engines detect deadlocks and roll back transactions to break circular waits \citep{grechanik}.
Haskell provides exceptions for certain permanently blocked computations; for example, a thread blocked on an \emph{MVar} can receive an exception when no other references to that variable remain.\footnote{\href{https://downloads.haskell.org/~ghc/7.10.3/docs/html/libraries/base-4.8.2.0/Control-Exception.html}{GHC 7.10.3 exception documentation.}}
Go's runtime can report a global deadlock as a fatal error, rather than a recoverable exception.\footnote{\href{https://go.dev/src/runtime/proc.go}{Go runtime deadlock detection.}}

Deadlock exceptions make an otherwise hidden failure visible and can support application-specific handling when recovery is appropriate.
Separating detection from automatic recovery avoids requiring a general rollback mechanism, which may be costly, inapplicable to some side effects, or conceal bugs that need correction.
Reliable reporting is especially valuable in non-interactive systems, where a deadlock may disable one behavior without stopping the entire application.
\citet{ian} describes such a case at CloudFlare: a two-thread deadlock prevented an expected behavior and was difficult to diagnose because the application continued running.

This thesis investigates deadlock exceptions as a way to help developers identify and fix concurrency bugs.
The implementation targets Java, but the motivation extends to other languages.
